% ID: FIS-URT-ANE-002
% Titolo: Inseguimento su piano inclinato e urto anelastico
% Disciplina: Fisica
% Area: Quantita di moto e urti
% Argomento: Urti anelastici
% Sottoargomento: Moto accelerato con attrito seguito da urto completamente anelastico
% Classe: Terza liceo scientifico
% Difficolta: 4
% Tipo: problema_numerico
% Risultato: t = 1,08 s; V = 2,56 m/s
% Tempo_stimato: 13 min
% Competenze: dinamica sul piano inclinato, cinematica relativa, conservazione della quantita di moto
% Prerequisiti: attrito dinamico, piano inclinato, moto uniformemente accelerato, urto completamente anelastico
% Tag: fisica, quantita_moto, urti_anelastici, piano_inclinato, attrito, inseguimento
% Fonte: rielaborazione_da_traccia_fornita_dal_docente
% Autore: Riccardo Carli
% Licenza: CC-BY-SA-4.0

\begin{esercizio}
Su un piano inclinato di \(30^\circ\) si trovano due corpi \(A\) e \(B\),
entrambi di massa \(m=2{,}00\,\mathrm{kg}\), inizialmente fermi e separati da
una distanza \(d=1{,}00\,\mathrm{m}\) misurata lungo il piano. Il corpo \(A\)
si trova piu in alto di \(B\).

I coefficienti di attrito dinamico con il piano sono
\(\mu_A=0{,}200\) e \(\mu_B=0{,}400\). I due corpi vengono lasciati scivolare
simultaneamente. Poiche \(A\) accelera maggiormente, raggiunge \(B\) e lo urta
in modo completamente anelastico.

Determina:
\begin{enumerate}
\item il tempo trascorso prima dell'urto;
\item la velocita comune dei due corpi immediatamente dopo l'urto.
\end{enumerate}
\end{esercizio}

\begin{soluzione}
Assumiamo positivo il verso diretto verso il basso lungo il piano. Per un corpo
che scivola lungo un piano inclinato con attrito dinamico:
\[
a=g(\sin\theta-\mu\cos\theta).
\]
Con \(\theta=30^\circ\) si ottiene per \(A\):
\[
a_A=9{,}81\left(0{,}500-0{,}200\cdot0{,}866\right)
\simeq3{,}21\,\mathrm{m/s^2}.
\]
Per \(B\):
\[
a_B=9{,}81\left(0{,}500-0{,}400\cdot0{,}866\right)
\simeq1{,}51\,\mathrm{m/s^2}.
\]

Poniamo l'origine nella posizione iniziale di \(A\). Poiche entrambi partono da
fermi:
\[
x_A(t)=\frac12a_At^2,
\qquad
x_B(t)=d+\frac12a_Bt^2.
\]
L'urto avviene quando \(x_A=x_B\):
\[
\frac12(a_A-a_B)t^2=d.
\]
Quindi
\[
t=\sqrt{\frac{2d}{a_A-a_B}}
=\sqrt{\frac{2\cdot1{,}00}{3{,}21-1{,}51}}
\simeq1{,}08\,\mathrm{s}.
\]

Subito prima dell'urto le velocita valgono
\[
v_A=a_At\simeq3{,}48\,\mathrm{m/s},
\qquad
v_B=a_Bt\simeq1{,}64\,\mathrm{m/s}.
\]

L'urto e completamente anelastico, quindi i due corpi proseguono insieme con
velocita \(V\). Durante l'urto si conserva la quantita di moto lungo il piano:
\[
mv_A+mv_B=(2m)V.
\]
Poiche le masse sono uguali:
\[
V=\frac{v_A+v_B}{2}
=\frac{3{,}48+1{,}64}{2}
\simeq2{,}56\,\mathrm{m/s}.
\]
La velocita finale e diretta verso il basso lungo il piano inclinato.
\end{soluzione}
